% 第08b节：黑洞物理 - 信息悖论、蒸发和内部结构
% UFT对黑洞谜题的解决
\section{黑洞物理：信息、蒸发和内部结构}
\label{sec:bh_physics_complete}

\subsection{引言}

黑洞对量子引力提出最严峻的挑战：
\begin{itemize}
    \item 信息悖论：落入黑洞的量子信息会发生什么？
    \item 奇点问题：物理在$r = 0$处崩溃吗？
    \item 蒸发终点：黑洞蒸发的最终状态是什么？
\end{itemize}

UFT通过10维内空间几何解决所有三个问题。

\subsection{信息悖论的解决}

\subsubsection{标准悖论}

在半经典引力中，霍金辐射是热辐射，暗示：
\begin{itemize}
    \item 纯态$\to$混合态演化（违反量子力学）
    \item 信息被毁灭（违反幺正性）
    \item Page曲线被违反
\end{itemize}

\subsubsection{UFT解决：信息守恒}

\begin{insight}[信息流向内空间]
当物质落入黑洞时，它并未消失——而是过渡到10维内空间。信息被保留但对4D观察者不可达。
\end{insight}

\begin{theorem}[UFT中的幺正演化]
包括4D互反空间和10D内空间的总演化是幺正的：
\begin{equation}
    S_{\text{total}} = S_{4D} \otimes S_{10D}, \quad S_{\text{total}}^\dagger S_{\text{total}} = \mathbb{1}
\end{equation}
\end{theorem}

\begin{proof}
完整量子态是$\mathcal{H}_{4D} \otimes \mathcal{H}_{10D}$中的矢量。黑洞形成和蒸发对应于该积空间中的幺正变换：
\begin{equation}
    |\Psi_{\text{in}}\rangle_{4D} \otimes |0\rangle_{10D} \to |\text{BH}\rangle_{4D} \otimes |\Psi\rangle_{10D} \to |\Psi_{\text{out}}\rangle_{4D} \otimes |0\rangle_{10D}
\end{equation}

对4D观察者看来"丢失"的信息实际上存储在10维内空间中。当黑洞蒸发时，内空间信息原则上可以从霍金辐射关联中重建。
\end{proof}

\subsubsection{Page曲线}

UFT中霍金辐射的纠缠熵遵循Page曲线：
\begin{equation}
    S_{\text{rad}}(t) = \min\left[S_{\text{BH}}(0), S_{\text{BH}}(t)\right] - \Delta S_{\text{corr}}
\end{equation}

其中$\Delta S_{\text{corr}} \sim k_B \ln d_{\text{int}}$计入内空间关联。

\subsection{奇点的解决}

\subsubsection{经典奇点}

在GR中，Schwarzschild度规在$r = 0$处有曲率奇点：
\begin{equation}
    R_{\mu\nu\rho\sigma}R^{\mu\nu\rho\sigma} \sim \frac{12r_s^2}{r^6} \to \infty \quad \text{当} \quad r \to 0
\end{equation}

\subsubsection{UFT解决：向内分形展开}

\begin{insight}[奇点是视角效应]
$r = 0$处的"奇点"不是无限密度的物理点。它是4D观察者对内空间展开以容纳能量的区域的视角。
\end{insight}

\begin{theorem}[UFT黑洞的内部结构]
视界内部（$r \u003c r_s$），度规被扭转修正：
\begin{equation}
    ds^2 = -\left(1 - \frac{r_s}{r}\right)dt^2 + \left(1 - \frac{r_s}{r} + \tau_0^2\frac{r_s^2}{r^2}\right)^{-1}dr^2 + r^2 d\Omega^2
\end{equation}

奇点被规则区域取代，其中$r_{\text{int}} \to \infty$当$r \to 0$。
\end{theorem}

\begin{proof}
扭转场贡献有效压力阻止坍缩：
\begin{equation}
    P_{\text{eff}} = P_{\text{matter}} + P_{\tau} = P_{\text{matter}} + \frac{\tau_0^2}{\kappa^2 r^4}
\end{equation}

在$r \sim r_s \tau_0$，扭转压力平衡引力吸引，创建规则的"反弹"区域而非奇点。
\end{proof}

\subsubsection{防火墙悖论的解决}

"防火墙"悖论暗示落入观察者遇到视界处的高能辐射。在UFT中：
\begin{itemize}
    \item 没有防火墙——视界是平滑的几何过渡
    \item 落入观察者平滑地过渡到内空间
    \item 等效原理被保留
\end{itemize}

\subsection{黑洞蒸发}

\subsubsection{修正的霍金温度}

UFT中的霍金温度被扭转场修正：
\begin{equation}
    T_H = \frac{\hbar c^3}{8\pi G M} \cdot \left(1 + \tau_0^2 \cdot \frac{M_P^2}{M^2}\right)
\end{equation}

对于天体物理黑洞（$M \gg M_P$），修正可忽略。对于原初黑洞（$M \sim 10^{15}$ g），修正可能显著。

\subsubsection{蒸发率}

质量损失率为：
\begin{equation}
    \frac{dM}{dt} = -\frac{\hbar c^4}{G^2 M^2} \cdot g_* \cdot f_\tau(M/M_P)
\end{equation}

其中$g_*$是自由度数，$f_\tau$是扭转抑制因子。

\subsubsection{蒸发的终点}

\begin{theorem}[黑洞终点是遗迹或爆发]
当$M \to M_{\text{crit}} \sim M_P/\tau_0$，蒸发行为改变：
\begin{itemize}
    \item 如果$\tau_0 \ll 1$：黑洞在最终爆发中完全蒸发（$\gamma$射线）
    \item 爆发能量：$E_{\text{burst}} \sim M_{\text{crit}} c^2 \sim 10^{16}$ GeV
\end{itemize}
\end{theorem}

\subsubsection{蒸发的观测信号}

\begin{enumerate}
    \item \textbf{$\gamma$射线暴}：$M \sim 10^{15}$ g PBH的最终爆发
    \item \textbf{宇宙射线}：蒸发PBH的高能粒子
    \item \textbf{中微子流}：霍金辐射的特征谱
\end{enumerate}

\subsection{黑洞内部}

\subsubsection{几何结构}

UFT黑洞内部有分层结构：

\begin{enumerate}
    \item \textbf{外区}（$r_s \u003e r \u003e r_s\tau_0$）：经典GR近似有效
    \item \textbf{过渡区}（$r \sim r_s\tau_0$）：扭转效应变得重要
    \item \textbf{内区}（$r \u003c r_s\tau_0$）：10D几何主导
    \item \textbf{核心}（$r \to 0$）：平滑、非奇异区域，$r_{\text{int}} \to \infty$
\end{enumerate}

\subsubsection{Penrose图}

UFT黑洞的因果结构被修正：
\begin{itemize}
    \item 奇点被类时边界（"内无穷"）取代
    \item 信息原则上可从内空间逃逸回4D
    \item 整体结构类似于连接渐近区域的"虫洞"
\end{itemize}

\subsection{旋转和带电黑洞}

\subsubsection{Kerr-Newman修正}

对于旋转、带电黑洞，UFT度规推广为：
\begin{equation}
    ds^2 = -\frac{\Delta_\tau}{\rho^2}\left(dt - a\sin^2\theta \, d\phi\right)^2 + \frac{\rho^2}{\Delta_\tau}dr^2 + \rho^2 d\theta^2 + \frac{\sin^2\theta}{\rho^2}\left(a \, dt - (r^2+a^2)d\phi\right)^2
\end{equation}

其中：
\begin{equation}
    \Delta_\tau = r^2 - 2Mr + a^2 + Q^2 + \tau_0^2 r_s^2
\end{equation}

\subsubsection{内Cauchy视界}

内视界（$r_-$）被扭转场稳定：
\begin{equation}
    r_{\pm} = M \pm \sqrt{M^2 - a^2 - Q^2 - \tau_0^2 r_s^2}
\end{equation}

质量膨胀不稳定性被$\tau_0^2$项抑制。

\subsection{量子方面}

\subsubsection{Bekenstein-Hawking熵}

UFT中的黑洞熵包括视界面积和内空间体积的贡献：
\begin{equation}
    S_{\text{BH}} = \frac{A}{4G} + S_{\text{int}} = \frac{\pi r_s^2}{G} + k_B \ln \mathcal{N}_{\text{int}}
\end{equation}

其中$\mathcal{N}_{\text{int}}$是内空间微观态数。

\subsubsection{拉伸视界}

有效的"拉伸视界"位于：
\begin{equation}
    r_{\text{stretched}} = r_s + \ell_P^2/r_s + \tau_0 r_s
\end{equation}

这是量子效应变得重要的位置。
